%% chapters/theory/chiral_higher_deligne.tex
%%
%% The Chiral Higher Deligne theorem package.
%%
%% Lane discipline:
%% - cohomological E_2 on chiral Hochschild cochains in general;
%% - chain-level E_2 only on the associator-fixed formality locus;
%% - E_3 only after explicit topologization data, not bare L_0
%% narration;
%% - Theorem H remains the Vol I family-support theorem,
%% while the E_3-PBW reinterpretation remains an open route.
%%
%% The DS-Hochschild bridge closes the brace / E_2 compatibility gap
%% and inherits any topologized E_3 layer only in the settings where the
%% topologization datum is already proved.
%%
%% Keeps the native/derived E_n distinction: E_3 on
%% Z^{der}_{ch}(A), never on A); three Hochschilds (ChirHoch /
%% HH*_mode / H*_GF); algebraic/geometric SC^{ch,top} models
%% (End^{ch}_A formal-Laurent vs. FM_k(C)-integration); chiral vs.
%% topological Deligne-Tamarkin (associator disclosure mandatory).

%% Local macro safety: providecommand only.
\providecommand{\SCchtop}{\mathsf{SC}^{\mathrm{ch},\mathrm{top}}}
\providecommand{\SCchtopKos}{\mathsf{SC}^{\mathrm{ch},\mathrm{top},!}}
\providecommand{\Fact}{\mathrm{Fact}}
\providecommand{\CoFact}{\mathrm{CoFact}}
\providecommand{\FM}{\mathrm{FM}}
\providecommand{\Conf}{\mathrm{Conf}}
\providecommand{\Ran}{\mathrm{Ran}}
\providecommand{\Obs}{\mathsf{Obs}}
\providecommand{\Hoch}{\mathrm{Hoch}}
\providecommand{\ChirHoch}{\mathrm{ChirHoch}}
\providecommand{\Bulk}{\mathrm{Bulk}}
\providecommand{\Bdy}{\mathrm{Bdy}}
\providecommand{\MC}{\mathrm{MC}}
\providecommand{\QME}{\mathrm{QME}}
\providecommand{\End}{\mathrm{End}}
\providecommand{\Mod}{\mathrm{Mod}}
\providecommand{\id}{\mathrm{id}}
\providecommand{\cA}{\mathcal{A}}
\providecommand{\cB}{\mathcal{B}}
\providecommand{\cC}{\mathcal{C}}
\providecommand{\cE}{\mathcal{E}}
\providecommand{\cF}{\mathcal{F}}
\providecommand{\cO}{\mathcal{O}}
\providecommand{\cZ}{\mathcal{Z}}
\providecommand{\cM}{\mathcal{M}}
\providecommand{\fg}{\mathfrak{g}}
\providecommand{\Lie}{\mathsf{Lie}}
\providecommand{\Ass}{\mathsf{Ass}}
\providecommand{\Com}{\mathsf{Com}}
\providecommand{\Free}{\mathrm{Free}}
\providecommand{\RHom}{\mathbf{R}\mathrm{Hom}}
\providecommand{\Barch}{\mathrm{B}^{\mathrm{ch}}}
\providecommand{\Cobarch}{\Omega^{\mathrm{ch}}}
\providecommand{\ChirAss}{\mathsf{ChirAss}}

\chapter{The Chiral Higher Deligne Conjecture}
\label{ch:chiral-higher-deligne}

%%%%%%%%%%%%%%%%%%%%%%%%%%%%%%%%%%%%%%%%%%%%%%%%%%%%%%%%%%%%%%%%%%%%%%%%%%%%%
\section*{Prologue}
%%%%%%%%%%%%%%%%%%%%%%%%%%%%%%%%%%%%%%%%%%%%%%%%%%%%%%%%%%%%%%%%%%%%%%%%%%%%%

The classical Deligne conjecture delivers an $E_2$-algebra structure
on Hochschild cochains up to associator (Tamarkin, Kontsevich,
McClure--Smith~\cite{McClureSmith02}, Berger--Fresse, Hu--Kriz--Voronov), one dimension
short of holomorphic factorisation on a curve. An $E_1$-chiral
algebra $\cA$ on a smooth complex curve $X$ has a derived chiral
centre
$Z^{\mathrm{der}}_{\mathrm{ch}}(\cA) := \ChirHoch^\bullet(\cA, \cA)$
(Theorem~H, with support computed from $H_H(\cA;S)$). The bulk structure of
the universal holography part (cf.~\ref{part:holography})
separates into three inputs: the $E_2$-chiral brace action; the
two-coloured $\SCchtop$ bulk-boundary lift on
$(Z^{\mathrm{der}}_{\mathrm{ch}}(\cA),\cA)$ under the
two-coloured cobar-homotopy hypothesis; the strict
$E_3^{\mathrm{top}}$ refinement supplied by explicit topologisation
data. The class-$M$ global-triangle gap is exactly the failure of
these three inputs to collapse to a single theorem. The chapter
constructs the $E_2$-chiral bridge, isolates the strict $E_3$
chain-lift obligation, and treats family support as the output of a
separate chiral-$E_3$-PBW theorem for the derived centre.

The mixed chiral-topological action is constructed via pentagon
edges $(3)\leftrightarrow(4)\leftrightarrow(5)$ of the
$\SCchtop$-heptagon (Chapter~\ref{ch:sc-chtop-heptagon}). The
algebraic model $\ChirHoch^\bullet(\cA,\cA) =
\RHom_{\End^{\mathrm{ch}}_{\cA}}(\cA,\cA)$ has spectral parameters
in the formal-Laurent completion at the diagonal; the geometric
model via $\FM_k(\C)$-integration is connected by
Prop.~\ref{prop:chd-models-equivalent}. Multi-input brace
operations $B^{(n_1,\ldots,n_k)}$ realise chiral Deligne--Tamarkin
compatibility with the $\SCchtop$-pentagon edge
$(3)\leftrightarrow(4)$. The main theorem: the pair
$(Z^{\mathrm{der}}_{\mathrm{ch}}(\cA),\cA)$ carries the mixed
chiral-topological $\SCchtop$ datum through the heptagon under
the two-coloured cobar-homotopy hypothesis.  The $E_3$-PBW route
recovers a prescribed support set once the derived centre admits a
chiral-$E_3$-PBW presentation whose first page has that support.
Applications: the DS--Hochschild compatibility bridge at chain
level and the ambient-qualified universal-holography statement
for class~M.

Associator discipline: the chain-level $E_3$-action depends on a
choice of Drinfeld associator; the cohomological action is
associator-free (Remark~\ref{rem:chd-associator-discipline}).

%%%%%%%%%%%%%%%%%%%%%%%%%%%%%%%%%%%%%%%%%%%%%%%%%%%%%%%%%%%%%%%%%%%%%%%%%%%%%
\section{Algebraic and geometric models for $\ChirHoch^\bullet$}
\label{sec:chd-setup}
%%%%%%%%%%%%%%%%%%%%%%%%%%%%%%%%%%%%%%%%%%%%%%%%%%%%%%%%%%%%%%%%%%%%%%%%%%%%%

Two inequivalent-looking presentations of chiral Hochschild
cohomology coexist in the literature and earlier in this volume.
Both are correct; each is adapted to a different geometric
question. The two are fixed and bridged below.

\begin{definition}[Algebraic chiral Hochschild complex]%
\label{def:chd-alg}%
\index{chiral Hochschild cohomology!algebraic model}%
Let $\cA$ be an $E_1$-chiral algebra on a smooth curve $X$ with
structure morphisms $\mu_z : \cA \boxtimes \cA \to \Delta_! \cA$
parametrised by the punctured formal neighbourhood of the diagonal
$\Delta \subset X\times X$. The \emph{chiral endomorphism operad}
$\End^{\mathrm{ch}}_{\cA}$ is the one-coloured symmetric operad in
the ambient category $\mathrm{D}\text{-}\mathrm{mod}(X^\bullet)$ of
right D-modules on the finite powers of $X$: its $k$-ary component
is the ind-object
\[
\End^{\mathrm{ch}}_{\cA}[k] \;:=\;
\varinjlim_{n}\,
\Hom_{\mathrm{D}(X^k)}\bigl(\cA^{\boxtimes k},\,
\Delta^{(k)}_! \cA\,(n\!\cdot\!\partial\Delta^{(k)})\bigr),
\]
the multi-linear operations
\[
\phi_k : \cA^{\boxtimes k} \longrightarrow \Delta^{(k)}_! \cA
\]
with poles of bounded order along the diagonal divisor, i.e.\
polynomial in the formal-Laurent disc at the small diagonal
$\Delta^{(k)} \subset X^k$; operadic composition is composition of
such operations followed by the diagonal pushforward, and the
$\Sigma_k$-action permutes the factors of $X^k$. The spectral parameters
$z_1, \ldots, z_k \in X$ enter as formal variables in this
Laurent completion: the spectral parameters are formal-algebraic
variables in $\End^{\mathrm{ch}}_\cA$, not $\FM_k(\C)$ coordinates.
The local-global comparison making the parameters geometric is
Proposition~\ref{prop:chd-models-equivalent}.
The algebraic chiral Hochschild complex is
\begin{equation}\label{eq:chd-alg-def}
 C^\bullet_{\mathrm{ch,alg}}(\cA, \cA) \;:=\;
 \RHom_{\End^{\mathrm{ch}}_\cA}(\cA, \cA),
\end{equation}
with differential the sum of internal $\cA$-differential and
$\ChirAss$-cobar coderivation.
\end{definition}

\begin{definition}[Geometric chiral Hochschild complex]%
\label{def:chd-geom}%
\index{chiral Hochschild cohomology!geometric model}%
With notation as above, the \emph{geometric} chiral Hochschild complex
is the totalisation
\begin{equation}\label{eq:chd-geom-def}
 C^\bullet_{\mathrm{ch,geom}}(\cA, \cA) \;:=\;
 \operatorname{Tot} \Bigl[\,
 \bigoplus_{k \ge 0}
 \Gamma\bigl(\FM_k(\C),\,
 \omega_{\FM_k(\C)}\otimes
 \underline{\Hom}(\cA^{\boxtimes k}, \cA)\bigr)\,\Bigr],
\end{equation}
where $\FM_k(\C)$ is the Fulton-MacPherson compactification of the
ordered configuration space $\Conf_k(\C)$ and $\omega_{\FM_k(\C)}$ the
sheaf of top logarithmic forms. The typing of the coefficient sheaf
requires a choice: $\FM_k(\C)$ parametrises configurations in a fixed
affine coordinate, not points of $X^k$, so
$\underline{\Hom}(\cA^{\boxtimes k}, \cA)$ denotes the pullback of the
$k$-ary D-module component of Definition~\ref{def:chd-alg} along a
fixed \'etale coordinate $t\colon U \to \C$ on a chart $U \subseteq X$,
through the induced map $\Conf_k(U) \to \Conf_k(\C)$ and its extension
to the compactification. The geometric complex is defined
chart-locally; independence of the chart up to the transition action
of formal coordinate changes ($\mathrm{Aut}\,\hat{\mathcal O} = W_+$-torsor
descent) is part of the comparison hypothesis of
Proposition~\ref{prop:chd-models-equivalent}, not of this definition.
\end{definition}

\begin{proposition}[Equivalence of algebraic and geometric models;
\ClaimStatusConditional{} on the typing hypothesis
\textup{(H$_{\mathrm{typ}}$)}]%
\label{prop:chd-models-equivalent}%
\index{chiral Hochschild cohomology!algebraic = geometric}%
Assume the typing hypothesis \textup{(H$_{\mathrm{typ}}$)}: a fixed
\'etale coordinate $t\colon U \to \C$ on a chart $U \subseteq X$
identifying the $\FM_k(\C)$ configuration coordinates with the
diagonal Laurent variables $z_1,\ldots,z_k$ of
Definition~\ref{def:chd-alg}, together with descent data for the
$\mathrm{Aut}\,\hat{\mathcal O}$-action on the coefficient sheaf
making the chart-local geometric complex
\eqref{eq:chd-geom-def} glue over $X$. Under
\textup{(H$_{\mathrm{typ}}$)} the two models are canonically
quasi-isomorphic:
\[
 \iota_{\mathrm{a,g}}:
 C^\bullet_{\mathrm{ch,alg}}(\cA, \cA) \;\simeq\;
 C^\bullet_{\mathrm{ch,geom}}(\cA, \cA),
\]
and the equivalence intertwines the $\ChirAss$-cobar differential with
the Stokes differential on $\FM_\bullet(\C)$.
\end{proposition}

\begin{proof}
Four steps. Step 1: identification of $k$-ary operations. A
$\Hom$-value on $\cA^{\boxtimes k} \to \cA$ that is polynomial in the
formal-Laurent completion at $\Delta^{(k)}$ extends uniquely to a
section of $\omega_{\FM_k(\C)}\otimes \underline{\Hom}(\cA^{\boxtimes k}, \cA)$
across the compactification divisor: the divisor components of
$\FM_k(\C)$ are indexed by rooted trees, and the formal-Laurent expansion at
each tree stratum realises the corresponding nested collision limit.
This is the chiral analogue of Getzler--Jones for $E_n$
(cf.~Kontsevich 2006, arXiv:math/0608180, \S4, for rational
weights; rational weights are replaced by logarithmic de Rham
periods).

Step 2: differential compatibility. The $\ChirAss$-cobar coderivation
$d_{\mathrm{cobar}}$ on~\eqref{eq:chd-alg-def} sums over placements of
$\mu_z$ into a given $\phi_k$. On the geometric side, the Stokes
differential $d_{\mathrm{Stokes}}$ sums over codimension-one strata of
$\FM_k(\C)$, each corresponding to a nested collision. Bijection
between placements and strata yields
$\iota_{\mathrm{a,g}}\circ d_{\mathrm{cobar}} = d_{\mathrm{Stokes}}\circ \iota_{\mathrm{a,g}}$.

Step 3: quasi-isomorphism on $E_1$ pages. Both complexes carry the
weight filtration by number of inputs $k$; the associated graded on
either side is $\bigoplus_k \underline{\Hom}(\cA^{\boxtimes k}, \cA)$
with the internal $\cA$-differential. The identity map on $E_1$ pages
is a quasi-isomorphism.

Step 4: convergence. Both filtrations are bounded below by $k = 0$
and exhaustive; the $E_1$ quasi-isomorphism lifts to a quasi-isomorphism
of totalisations.
\end{proof}

\begin{convention}[Dealiasing]%
\label{conv:chd-dealias}%
In the remainder of this chapter $\ChirHoch^\bullet(\cA, \cA)$
denotes either model interchangeably, bridged by
Prop.~\ref{prop:chd-models-equivalent}. The notation
$C^\bullet_{\mathrm{ch,alg}}$ is used when the formal-Laurent /
operadic structure is in play; $C^\bullet_{\mathrm{ch,geom}}$ when
divisor / Stokes / integration arguments are required.
\end{convention}

\begin{remark}[Distinguish three Hochschilds]%
\label{rem:chd-three-hochschilds}%
\index{Hochschild!three theories}%
The complex $\ChirHoch^\bullet(\cA, \cA)$ is the
\emph{chiral} Hochschild complex of the $E_1$-chiral algebra $\cA$ on
the curve $X$. It is to be distinguished from
\begin{itemize}
\item $HH^\bullet_{\mathrm{mode}}(\cA_0, \cA_0)$, the \emph{mode-algebra}
Hochschild complex of the associative mode algebra $\cA_0$ underlying a
vertex algebra ($\cA_0 = \bigoplus_n \cA_{(n)}$ with its Cauchy product;
this is the topological Hochschild of an $E_1$-topological algebra);
\item $H^\bullet_{\mathrm{GF}}(\cA)$, the Gel'fand-Fuchs continuous
cohomology of the formal vector-field algebra, invoked on the small
disc $D \subset X$ when $\cA$ is the Virasoro or $\cW$-algebra at
generic central charge.
\end{itemize}
The three agree on their intersection of validity only after the
appropriate spectral / derived completion, and are generically
inequivalent: $HH^\bullet_{\mathrm{mode}}$ knows the associative
structure at the level of modes but is blind to OPE spectral data;
$H^\bullet_{\mathrm{GF}}$ is unbounded above (Fuks 1986).  A family
datum $H_H(\cA;S)$ identifies $\ChirHoch^\bullet(\cA)$ with a model
supported in~$S$ (Theorem~H). The Weyl algebra identity $HH^\bullet(\mathrm{Weyl}_n)
= \Bbbk$ concentrated in degree $0$ (Fresse-Tsygan) is a statement
about the MODE algebra; the corresponding chiral statement is
$\ChirHoch^\bullet(\mathrm{Weyl}^{\mathrm{ch}}_n) = \Bbbk$ in degree
$0$ and a $g$-dimensional piece in degree $1$
(Vol.~I, Prop.~\ref*{V1-prop:chirhoch1-affine-km} specialised to the
Heisenberg limit). The \emph{unbounded} object occurring in physics is
Gel'fand-Fuchs, NOT THH: the latter has no analogue of Theorem~H
because THH is not a chiral invariant at all.
\end{remark}

%%%%%%%%%%%%%%%%%%%%%%%%%%%%%%%%%%%%%%%%%%%%%%%%%%%%%%%%%%%%%%%%%%%%%%%%%%%%%
\section{Brace operations and chiral Deligne-Tamarkin compatibility}
\label{sec:chd-braces}
%%%%%%%%%%%%%%%%%%%%%%%%%%%%%%%%%%%%%%%%%%%%%%%%%%%%%%%%%%%%%%%%%%%%%%%%%%%%%

The $E_3$ action arises from a chiral-operadic refinement of
Kontsevich's brace operations. The braces are constructed as
explicit chain-level maps; Stasheff coherences hold through
degree~$4$ and compare to the pentagon edge
$(3)\leftrightarrow(4)$ of the $\SCchtop$-heptagon.

\begin{definition}[Multi-input chiral brace operations]%
\label{def:chd-braces}%
\index{brace operations!chiral multi-input}%
For integers $k \ge 1$ and $n_1, \ldots, n_k \ge 0$ with
$p_0 = n_1 + \cdots + n_k$, the multi-input brace operation
\begin{equation}\label{eq:chd-braces}
 B^{(n_1, \ldots, n_k)}:
 \ChirHoch^{p_0+k}(\cA,\cA) \otimes
 \ChirHoch^{p_1}(\cA,\cA) \otimes \cdots \otimes
 \ChirHoch^{p_k}(\cA,\cA)
 \;\longrightarrow\;
 \ChirHoch^{p_0+p_1+\cdots+p_k+1-k}(\cA,\cA)
\end{equation}
is the chain map defined as follows. Represent the left operand as a
$(p_0+k)$-ary chiral operation $\phi$ in
$\End^{\mathrm{ch}}_\cA[p_0+k]$; represent the right operands as
$p_i$-ary chiral operations $\psi_i$. Choose $k$ disjoint windows in
the $p_0+k$ input slots of $\phi$, each window of size $n_i$ followed
by one insertion slot. The brace value
$\phi\{\psi_1, \ldots, \psi_k\}_{(n_1, \ldots, n_k)}$ is the chiral
operation obtained by inserting $\psi_i$ in the $i$-th insertion slot
and leaving $n_i$ consecutive original slots above each insertion
untouched, summed over the chosen window positions with signs
$(-1)^{\sum p_i n_j[j>i]}$.

In the geometric model (Def.~\ref{def:chd-geom}), \eqref{eq:chd-braces}
is the pullback of the product of forms along the diagonal nesting
$\FM_k(\C)\hookrightarrow \FM_{p_0+k}(\C)$ with the $\psi_i$ inserted
at the $k$ nested points.
\end{definition}

\begin{proposition}[Chain map with explicit associator disclosure]%
\label{prop:chd-brace-chain-map}%
\index{brace operations!chain map}%
For every fixed $(n_1, \ldots, n_k)$ the operation
$B^{(n_1, \ldots, n_k)}$ commutes with the chiral Hochschild
differential $d_{\mathrm{ch}} = d_\cA + d_{\mathrm{cobar}}$:
\[
 d_{\mathrm{ch}}\bigl(\phi\{\psi_1, \ldots, \psi_k\}\bigr) \;=\;
 (d_{\mathrm{ch}}\phi)\{\psi_1, \ldots, \psi_k\} \;+\;
 \sum_{i=1}^k (-1)^{\epsilon_i}\, \phi\{\psi_1, \ldots, d_{\mathrm{ch}}\psi_i, \ldots, \psi_k\},
\]
with explicit Koszul sign $\epsilon_i = |\phi| + \sum_{j<i} |\psi_j|$.
\end{proposition}

\begin{proof}
$d_\cA$-compatibility is immediate from its being an internal
differential commuting with insertion. $d_{\mathrm{cobar}}$-compatibility
follows from the chiral Leibniz identity for $\mu_z$ (OPE
associativity of $\cA$): the cobar differential sums over placements
of $\mu_z$ into a given operation, and a placement either lands on a
non-window slot of $\phi$ (producing
$(d_{\mathrm{cobar}}\phi)\{\psi\}$), on an insertion slot
(producing $\phi\{d_{\mathrm{cobar}}\psi_i\}$), or across the window
boundary. Across-boundary placements cancel in pairs by the chiral
Leibniz identity for $\mu_z$ (equivalently, OPE associativity of $\cA$
expressed as a constraint on pair placements of the chiral
coproduct structure, universal for any $\ChirAss$-algebra).
\end{proof}

\begin{proposition}[Stasheff coherence at degree 4]%
\label{prop:chd-stasheff-4}%
\index{brace operations!Stasheff coherence}%
The multi-input braces satisfy the Stasheff identities up to degree $4$:
\begin{align}
 \phi\{\psi\}\{\chi\} \;-\; \phi\{\psi\{\chi\}\} \;-\;
 \phi\{\psi,\chi\} \;-\; (-1)^{|\psi||\chi|}\phi\{\chi,\psi\}
 \;&=\; (d_{\mathrm{ch}}\text{-exact}), \\
 \bigl(\phi\{\psi_1,\psi_2\}\bigr)\{\chi\} \;-\;
 \phi\{\psi_1\{\chi\},\psi_2\} \;-\;
 (-1)^{|\psi_1||\chi|} \phi\{\psi_1,\psi_2\{\chi\}\}
 \;-\; \phi\{\psi_1,\chi,\psi_2\} \;&=\; (d_{\mathrm{ch}}\text{-exact}).
\end{align}
\end{proposition}

\begin{proof}
Each coherence reduces, by Prop.~\ref{prop:chd-models-equivalent}, to
a Stokes identity on a three-dimensional face of
$\FM_{p+k+1}(\C)$. The face boundary decomposes into the four summand
terms, and Stokes' theorem provides the explicit chain homotopy as
the integral over the three-cell. The cell integrand is bounded
(logarithmic singularities cancel in $\omega_{\FM}$); hence the
integral is a legitimate chain-level $d_{\mathrm{ch}}$-boundary.
\end{proof}

\begin{remark}[Stasheff coherences beyond degree $4$ are open:
McClure--Smith combinatorics]%
\label{rem:chd-stasheff-higher-degree}%
\index{brace operations!Stasheff coherence!higher degree}%
Only the two brace identities of
Proposition~\ref{prop:chd-stasheff-4} are verified in this volume.
The degree-$n$ Stasheff coherence for $n \ge 5$ would reduce, via
Proposition~\ref{prop:chd-models-equivalent}, to a Stokes identity on
an $n$-dimensional face of $\FM_{m}(\C)$ for appropriate $m$; but the
required matching --- between the full tree of degree-$n$ Stasheff
summands and the codimension pattern of the Fulton--MacPherson
stratification, with compatible signs and cancelling logarithmic
singularities on every cell --- is exactly the sequence-operad
combinatorics that McClure--Smith carry out for the topological brace
operad \cite{McClureSmith02}, and that combinatorics has \emph{not}
been carried out here for the chiral braces; the corresponding
open-closed coherence structure at all arities is the
Kajiura--Stasheff open-closed homotopy algebra
\cite{KajiuraStasheff06}, likewise not constructed here in the
chiral setting. The extension to $n \ge 5$ is therefore an
open obligation, not a corollary of the degree-$4$ template; no
statement in this volume may cite this remark as licence for
``Stasheff coherences through all degrees.'' What is available
unconditionally is the degree-$\le 4$ package of
Proposition~\ref{prop:chd-stasheff-4} together with the
associator-fixed strictification of
Remark~\ref{rem:chd-brace-strict-coordinates}.
\end{remark}

\begin{remark}[Brace identities are chain-level, not strict;
promoted to convention]%
\label{rem:chd-brace-strict-coordinates}%
\index{brace operations!strict coordinate convention}%
The identities of Prop.~\ref{prop:chd-stasheff-4} hold up to
\emph{explicit} chain-level homotopies, not as strict equalities.
Whenever a subsequent statement uses strict braces, it passes through
this coherently chosen strictification. The strictification is the data
of a Drinfeld associator $\Phi \in \mathrm{GRT}_1(\Bbbk)$; different
associators produce different strict-brace coordinates, all
chain-equivalent. Cohomologically, all choices agree.
\end{remark}

\begin{theorem}[Chiral Deligne-Tamarkin compatibility]%
\label{thm:chd-deligne-tamarkin}%
\index{Deligne-Tamarkin!chiral}%
Fix a Drinfeld associator $\Phi$ and let $\alpha_\Phi$ be the
associated $A_\infty$-morphism from the little-disks operad $E_2$ to
the brace operad. The chiral braces of
Def.~\ref{def:chd-braces} realise the $\SCchtop$-Koszul-dual
cooperad's closed-colour component: the pentagon edge
$(3)\leftrightarrow(4)$ of the heptagon
(Chapter~\ref{ch:sc-chtop-heptagon}), which identifies the
$E_2$-chiral cooperad with the closed colour of $(\SCchtop)^!$,
becomes the identity
\[
 \SCchtop\text{-brace on } \ChirHoch^\bullet(\cA,\cA)
 \;=\; \alpha_\Phi(\text{Kontsevich brace on } \ChirHoch^\bullet_{\mathrm{form}}(\cA,\cA)),
\]
where $\ChirHoch^\bullet_{\mathrm{form}}$ is the formal-disc
specialisation. The right-hand side is the topological Deligne-Tamarkin
output at the formal disc; the left-hand side is the chiral refinement
globalised over $X$. The equality is an equality of
$A_\infty$-morphisms up to the associator-fixed strictification of
Remark~\ref{rem:chd-brace-strict-coordinates}.
\end{theorem}

\begin{proof}
Forward ($\subseteq$). The chiral brace is defined on the
Laurent-polynomial model and restricts at the formal disc to Kontsevich's
brace; Prop.~\ref{prop:chd-models-equivalent} restricted to a formal
disc $D \subset X$ yields the formal-disc specialisation on the nose.

Reverse ($\supseteq$). Given Kontsevich's brace on the formal disc,
its globalisation to $X$ is governed by the $\SCchtop$-cooperad
bar-duality; this is the content of the pentagon edge
$(3)\leftrightarrow(4)$, proved in Chapter~\ref{ch:sc-chtop-heptagon}
(Proposition~\ref*{prop:heptagon-edge-34}), using holomorphic descent
for the curve $X$ and the Stokes compatibility of
$\omega_{\FM}$. Applying $\alpha_\Phi$ strictifies the two sides into
the same $A_\infty$-morphism.

Chiral versus topological. The distinction is that the
chiral braces use OPE residues on $\FM_k(\C)$, not little $2$-disks on
$\R^2$; on cohomology both reduce to the $E_2$-Gerstenhaber structure
(formality) and agree; at chain level they differ by the explicit
Stokes-integral homotopies of Prop.~\ref{prop:chd-stasheff-4}.
\end{proof}

%%%%%%%%%%%%%%%%%%%%%%%%%%%%%%%%%%%%%%%%%%%%%%%%%%%%%%%%%%%%%%%%%%%%%%%%%%%%%
\section{The Chiral Higher Deligne Theorem}
\label{sec:chd-main}
%%%%%%%%%%%%%%%%%%%%%%%%%%%%%%%%%%%%%%%%%%%%%%%%%%%%%%%%%%%%%%%%%%%%%%%%%%%%%

Lift the $E_2$-chiral action of the braces
(Thm~\ref{thm:chd-deligne-tamarkin}) to the mixed
chiral-topological bulk-boundary datum governed by $\SCchtop$. The
heptagon identifies the formal-disc brace model, the ordered
bar-cobar model, and the two-coloured bulk-boundary model; it does not
turn the directed two-coloured operad $\SCchtop$ itself into a single-colour
$E_3$ operad. A strict $E_3$-topological structure appears only after
the chiral direction has been topologised by explicit inner-stress
data.

\begin{convention}[Mixed chiral-topological datum in this chapter]%
\label{conv:chd-e3-chiral-meaning}%
\index{mixed chiral-topological datum!convention}%
The phrase \emph{$E_3$-chiral} in this chapter denotes the mixed
chiral-topological datum carried by a factorisation algebra on
$\R\times X$: chiral holomorphic factorisation on $X$, topological
factorisation on $\R$, and the two-coloured $\SCchtop$ action on the
pair $(Z^{\mathrm{der}}_{\mathrm{ch}}(\cA),\cA)$. It is not a Dunn
identification of $\SCchtop$ with a single-colour $E_3$ operad. When
$\cA$ carries inner conformal-vector data at non-critical level, the
chiral direction can become $Q$-exact and the mixed datum topologises
to a genuine single-colour $E_3^{\mathrm{top}}$ structure; Dunn
additivity is invoked only after that topological replacement has been
made. Without such topologisation data, the output remains the
generic-case $\SCchtop$ bulk-boundary structure.
\end{convention}

\begin{theorem}[Chiral Higher Deligne: braces, conditional two-coloured lift,
and topologised $E_3$; \ClaimStatusConditional on clauses
\textup{(2)}--\textup{(3)}; licensing tags
$\alpha+\beta+\gamma+\delta$]%
\label{thm:chiral-higher-deligne}%
\index{Chiral Higher Deligne}%
\index{Theorem H!chiral Higher Deligne bridge}%
\noindent\textbf{Status.}\enspace Clause \textup{(1)} is the
associator-fixed chiral Deligne--Tamarkin brace theorem. Clause
\textup{(2)} is conditional on the signed two-coloured cobar
contraction of Proposition~\ref{prop:chd-two-coloured-cobar-obstruction}.
Clause \textup{(3)} is conditional on the topologisation datum of
Convention~\ref{conv:chd-e3-chiral-meaning}.

Let $\cA$ be an $E_1$-chiral algebra on a smooth complex curve $X$ with
$\cA$ in the Koszul locus (Vol.~I, Definition~\ref*{V1-def:koszul-locus}).
After fixing a Drinfeld associator $\Phi \in \mathrm{GRT}_1(\Bbbk)$:
\begin{enumerate}
\item The chiral Hochschild complex
$Z^{\mathrm{der}}_{\mathrm{ch}}(\cA) = \ChirHoch^\bullet(\cA,\cA)$
carries the canonical chiral brace, equivalently $E_2$-chiral,
action.
\item Assume the heptagon hypotheses and a chain-level contracting
homotopy for the two-coloured $(\SCchtop)^!$-cobar complex
(Remark~\ref{rem:chd-universality-hypotheses} and
Proposition~\ref{prop:chd-two-coloured-cobar-obstruction}). Then the pair
\((Z^{\mathrm{der}}_{\mathrm{ch}}(\cA),\cA)\) carries the
two-coloured $\SCchtop$ bulk-boundary datum, and for every
$\SCchtop$-open/closed pair $(B,\cA)$ with bulk $B$ acting on
$\cA$ there is a unique homotopy class of $\SCchtop$-compatible maps
$B\to Z^{\mathrm{der}}_{\mathrm{ch}}(\cA)$ intertwining the
bulk-to-boundary structure.
\item Assume explicit topologisation data in the sense of
Convention~\ref{conv:chd-e3-chiral-meaning}. Then the mixed datum
refines to a strict $E_3^{\mathrm{top}}$ action. At chain level this
action depends on the associator $\Phi$; on cohomology it is
associator-independent.
\end{enumerate}
\end{theorem}

\begin{remark}[Two-coloured cobar homotopy]%
\label{rem:chd-universality-hypotheses}%
\label{rem:chd-universality-conditional}%
\index{Chiral Higher Deligne!universality hypotheses}%
Clause (2) invokes the contracting homotopy of the $(\SCchtop)^!$-cobar
complex on the Koszul locus. Vol.~I
\texttt{thm:chiral-positselski-weight-completed}
(\texttt{chapters/theory/theorem\_B\_scope\_platonic.tex:248}) establishes
the weight-completed chiral Positselski co-contra correspondence on the
\emph{single-colour} $E_1^{\mathrm{ch}}$ cooperad (with
\texttt{thm:chiral-positselski-7-2} an alias inside the same proved body).
The
$(\SCchtop)^!$-formulation required for clause~(2) lifts that
single-colour statement to the \emph{two-coloured} cooperad setting;
the two-coloured lift is a genuinely different theorem, not yet
established at chain level in Vol.~I or here. The exact proof
obligation is the construction of a contracting homotopy
\[
 h_{\mathrm{oc}}\colon
 \Omega((\SCchtop)^!) \longrightarrow \Omega((\SCchtop)^!)[-1]
\]
compatible with the open-to-closed void and the closed-to-open mixed
shuffle operations. Definition~\ref{def:chd-two-coloured-cobar-tree-homotopy}
spells out the signed tree formula; Proposition~\ref{prop:chd-two-coloured-cobar-obstruction}
records the exact obstruction class. On cohomology the $E_2$ brace
structure does not use this homotopy. The two-coloured universal
$\SCchtop$ statement and the strict $E_3$ refinement use precisely the
extra hypotheses stated in the theorem.
\end{remark}

\begin{definition}[Signed two-coloured cobar tree homotopy]%
\label{def:chd-two-coloured-cobar-tree-homotopy}%
\index{Boardman--Vogt!two-coloured cobar tree homotopy}%
Let
\[
 \mathcal C_{\mathrm{oc}}:=\Omega((\SCchtop)^!)
\]
be the free coloured cobar dg operad generated by the suspended
dual corollas of \((\SCchtop)^!\). For a homogeneous decorated
two-coloured rooted tree
\[
 T=(T;\{x_v\}_{v\in V(T)})
\]
write \(E_{\mathrm{int}}(T)=\{e_1<\cdots<e_N\}\) for the internal
edges ordered by the orientation line
\(\mathfrak{o}(T)=\det(\Bbbk^{E_{\mathrm{int}}(T)})\). If \(e=e_j\),
let \(T_e\) denote the tree obtained by inserting a
Boardman--Vogt edge-length parameter on \(e\) and replacing that edge
by the corresponding suspended cobar generator. The candidate
degree \(-1\) two-coloured homotopy is
\begin{equation}\label{eq:chd-hoc-tree-formula}
 h_{\mathrm{oc}}^{\mathrm{tree}}(T)
 =
 \sum_{e\in E_{\mathrm{int}}(T)}
 (-1)^{\sigma(e,T)}T_e,
 \qquad
 \sigma(e_j,T)
 =
 (j-1)+\sum_{v\prec e_j}|x_v|\pmod 2.
\end{equation}
Here \(v\prec e_j\) means that, in the chosen edge/vertex order,
the vertex decoration \(x_v\) crosses the suspended generator created
by expanding \(e_j\). The formula is set equal to zero on the
open-to-closed void:
\[
 h_{\mathrm{oc}}^{\mathrm{tree}}\bigl(
 \mathcal C_{\mathrm{oc}}(\mathsf{cl}^{p},\mathsf{op}^{q};\mathsf{cl})
 \bigr)=0\qquad(q>0).
\]
Changing the total order changes both \(\mathfrak{o}(T)\) and
\(\sigma(e,T)\) by the same Koszul sign, so the operator is an
orientation-line expression rather than an ordering-dependent
enumeration.
\end{definition}

\begin{proposition}[Obstruction measured by the tree homotopy]%
\label{prop:chd-two-coloured-cobar-obstruction}%
\index{Boardman--Vogt!two-coloured obstruction class}%
Let \(p\colon\mathcal C_{\mathrm{oc}}\to\mathcal C_{\mathrm{min}}\)
be the projection to the chosen minimal two-coloured cobar retract and
let \(i\colon\mathcal C_{\mathrm{min}}\to\mathcal C_{\mathrm{oc}}\)
be its inclusion, with \(p\circ i=\mathrm{id}_{\mathcal C_{\mathrm{min}}}\).
The chain-level two-coloured cobar-homotopy
hypothesis in Theorem~\ref{thm:chiral-higher-deligne}(2) is exactly
the assertion that \(h_{\mathrm{oc}}^{\mathrm{tree}}\) descends through
the Boardman--Vogt degeneracy relations and satisfies
\begin{equation}\label{eq:chd-hoc-contraction}
 d\,h_{\mathrm{oc}}^{\mathrm{tree}}
 +h_{\mathrm{oc}}^{\mathrm{tree}}d
 =
 \mathrm{id}-i\circ p,
\end{equation}
together with the strong-deformation-retract side conditions
\begin{equation}\label{eq:chd-hoc-side-conditions}
 p\,h_{\mathrm{oc}}^{\mathrm{tree}}=0,\qquad
 h_{\mathrm{oc}}^{\mathrm{tree}}\,i=0,\qquad
 \bigl(h_{\mathrm{oc}}^{\mathrm{tree}}\bigr)^2=0,
\end{equation}
together with the mixed-cooperadic derivation identity
\begin{equation}\label{eq:chd-hoc-delta-mix}
 \Delta_{\mathrm{mix}}h_{\mathrm{oc}}^{\mathrm{tree}}(x)
 =
 \sum h_{\mathrm{oc}}^{\mathrm{tree}}(x')\otimes x''
 +(-1)^{|x'|}x'\otimes h_{\mathrm{oc}}^{\mathrm{tree}}(x'')
 \quad
 \bigl(\Delta_{\mathrm{mix}}x=\sum x'\otimes x''\bigr).
\end{equation}
Its failure is the obstruction class
\begin{equation}\label{eq:chd-oc-obstruction-class}
 [o_{\mathrm{oc}}]\in
 H^{1}\operatorname{Cone}\!\left(
 \Omega((\SCchtop)^!)
 \longrightarrow
 \End(Z^{\mathrm{der}}_{\mathrm{ch}}(\cA),\cA)
 \right).
\end{equation}
The strict chain-level \(\SCchtop\) lift exists precisely after this
class is killed. Without that vanishing, the chapter has the
chain-level \(E_2\)-chiral brace action and the cohomological
two-coloured statement, not a strict chain-level Swiss-cheese
universal theorem.
\end{proposition}

\begin{proof}
The cubical differential on the Boardman--Vogt interval attached to
an internal edge has endpoints \(t_e=0\) and \(t_e=1\). In the cobar
model these are respectively the contracted-edge term and the
composed-vertex term. The sign in
\eqref{eq:chd-hoc-tree-formula} is the contraction sign for the
orientation line \(\mathfrak{o}(T)\), corrected by the Koszul sign
from moving the suspended cobar generator past the preceding vertex
decorations. Hence the boundary of a once-expanded edge gives the
summand of \(T-i\circ p(T)\) supported at that edge.

The low-arity calculation fixes the sign convention.  If
\(E_{\mathrm{int}}(T)=\{e_1\}\) and
\(a=\sum_{v\prec e_1}|x_v|\), then
\[
 h_{\mathrm{oc}}^{\mathrm{tree}}(T)=(-1)^aT_{e_1},
 \qquad
 d_{\mathrm{cube}}T_{e_1}=(-1)^a\bigl(T-i p(T)\bigr),
\]
where \(d_{\mathrm{cube}}\) is the interval differential, after
identifying the two endpoints of the Boardman--Vogt interval with the
unexpanded and contracted-edge representatives.  Thus the
one-edge part of
\(d h_{\mathrm{oc}}^{\mathrm{tree}}
+h_{\mathrm{oc}}^{\mathrm{tree}}d\)
is \(T-i p(T)\).  If
\(E_{\mathrm{int}}(T)=\{e_1<e_2\}\), the two interval-insertion terms
are
\[
 (-1)^{\sigma(e_1,T)+\sigma(e_2,T_{e_1})}T_{e_1,e_2}
 +
 (-1)^{\sigma(e_2,T)+\sigma(e_1,T_{e_2})}T_{e_2,e_1}.
\]
The canonical transposition
\(\mathfrak{o}(T_{e_1,e_2})\simeq-\mathfrak{o}(T_{e_2,e_1})\) gives
\[
 \sigma(e_1,T)+\sigma(e_2,T_{e_1})
 \equiv
 \sigma(e_2,T)+\sigma(e_1,T_{e_2})+1
 \pmod 2,
\]
so these two terms cancel.  The general cancellation is the same
two-edge calculation applied to every unordered pair of internal
edges; no new sign occurs in higher arity.

Terms with two expanded edges cancel in pairs: expanding \(e_i\) and
then \(e_j\) gives the opposite orientation-line sign from expanding
\(e_j\) and then \(e_i\). If a sub-tree has closed output and an open
input, the relevant arity of \(C_\bullet^{\log}\SCchtop\) is the zero
complex, so both sides of
\eqref{eq:chd-hoc-contraction} vanish there. Cutting a tree before or
after expanding an internal mixed edge gives exactly the two summands
in \eqref{eq:chd-hoc-delta-mix}; the displayed sign is the Koszul sign
for moving a degree \(-1\) homotopy past \(x'\). The normalized
Boardman--Vogt degeneracy quotient gives the side conditions
\eqref{eq:chd-hoc-side-conditions}: \(p\) kills an inserted interval
coordinate, \(i\) lands in the chosen minimal representatives, and a
second interval insertion is degenerate. Thus
\eqref{eq:chd-hoc-contraction}--\eqref{eq:chd-hoc-delta-mix} are
precisely the two-coloured cobar contraction identities. Applying the
convolution functor
\(\Hom(-,\End(Z^{\mathrm{der}}_{\mathrm{ch}}(\cA),\cA))\) identifies
their defect with the degree-one Maurer--Cartan obstruction class
\eqref{eq:chd-oc-obstruction-class}.
\end{proof}

\begin{proof}
Part (1). The $E_2$ brace action.

\emph{Stage 1: braces give $E_2$.} By Prop.~\ref{prop:chd-brace-chain-map},
Prop.~\ref{prop:chd-stasheff-4}, and Thm~\ref{thm:chd-deligne-tamarkin},
the multi-input braces assemble into an $E_2$-chiral action on
$\ChirHoch^\bullet(\cA, \cA)$ after associator-fixed strictification.

The strict $E_3$ refinement is not a consequence of the Koszul locus
alone. If the topologisation data of
Convention~\ref{conv:chd-e3-chiral-meaning} is supplied, then the
$E_2$ brace action combines with the named $E_1^{\mathrm{top}}$
direction by Dunn additivity and yields the strict
$E_3^{\mathrm{top}}$ action.

Part (2). Assume the obstruction class
\eqref{eq:chd-oc-obstruction-class} vanishes and choose the normalized
tree homotopy \(h_{\mathrm{oc}}^{\mathrm{tree}}\) satisfying
\eqref{eq:chd-hoc-contraction}--\eqref{eq:chd-hoc-delta-mix}. The
convolution complex controlling \(W(\SCchtop)\)-bulk-boundary maps
then retracts to the minimal two-coloured cobar model. Given
$(B, \cA)$, the $\SCchtop$-action of $B$ on
$\cA$ is given by multi-input mixed-colour brace operations
$\nu_k^{(m)}: B^{\otimes k}\otimes \cA^{\otimes m}\to \cA$. The adjoint
\[
 \nu_k^{(m)} \;\mapsto\; \hat{\nu}_k:
 B^{\otimes k} \to \underline{\Hom}(\cA^{\otimes m}, \cA)^{[m]} \to
 \ChirHoch^m(\cA, \cA)
\]
defines an $\SCchtop$-compatible map $B \to Z^{\mathrm{der}}_{\mathrm{ch}}(\cA)$ by
bar-Koszul duality for $\SCchtop$ (note that
$(\SCchtop)^! \ne \SCchtop$; we use $(\SCchtop)^!$ for adjunction,
and the dual is $(\Lie, \Ass)$ with shuffle-mixed closed colour). Two
lifts of the same bulk-boundary action produce chain-homotopic maps
because the difference is a boundary in this convolution retract,
with homotopy supplied by \(h_{\mathrm{oc}}^{\mathrm{tree}}\). Vol.~I
supplies the single-colour comparison; in this conditional theorem the
extra datum is exactly the two-coloured homotopy together with the
vanishing of \([o_{\mathrm{oc}}]\).

Part (3). The associator dependence is inherited from the
strictification of the braces (Remark~\ref{rem:chd-brace-strict-coordinates}).
Two associators $\Phi_1, \Phi_2 \in \mathrm{GRT}_1(\Bbbk)$ differ
by a Grothendieck-Teichm\"uller element
$\gamma = \Phi_2 \Phi_1^{-1}\in \mathrm{GRT}_1$; the induced
mixed chiral-topological actions differ by the $\gamma$-conjugation
automorphism, which is an identity on cohomology
(Drinfeld-Bar-Natan-Etingof) and chain-homotopic to the identity by
an explicit Berger-Fresse-type homotopy.
\end{proof}

\begin{remark}[Associator discipline]%
\label{rem:chd-associator-discipline}%
\index{associator!chiral Higher Deligne dependence}%
The cohomological derived centre is associator-free. The chain-level
derived centre depends on $\Phi$,
and the space of equivalences between chain-level centres is a
$\mathrm{GRT}_1$-torsor. Bar-side invariants ($\kappa$, shadow tower)
are associator-free. Cobar-side coproduct invariants
(Stasheff-$A_\infty$ coefficients of the brace coordinates) depend
on $\Phi$. Theorems of the form "$Z^{\mathrm{der}}_{\mathrm{ch}}(\cA)$
equals the bulk of theory $T$" must state the cohomological and
chain-level variants separately.
\end{remark}

\begin{remark}[Native versus derived mixed structure]%
\label{rem:chd-native-vs-derived}%
\index{E_n algebra!native vs derived}%
$\cA$ itself is $E_1$-chiral. The mixed chiral-topological
$\SCchtop$ structure lives on the pair
$(Z^{\mathrm{der}}_{\mathrm{ch}}(\cA),\cA)$ as an output of the
derived-centre construction. Writing ``$E_3$ on $\cA$'' is a type
error; writing ``mixed $\SCchtop$ bulk-boundary datum on
$(Z^{\mathrm{der}}_{\mathrm{ch}}(\cA),\cA)$'' is the correct
native/derived distinction. A single-colour $E_3^{\mathrm{top}}$
structure is a further topologised refinement, not the native output.
\end{remark}

% =====================================================================
%  Conditional topologisation theorem: the shared (T1)-(T5) proof
%  obligations underlying every concrete SCchtop -> E_3^top move
%  (Sugawara for class L, DS BRST for class L -> class M, KZ analytic
%  SDR for class M / PVA).  Makes residual (R3) of the typed master
%  functor of foundations.tex:thm:typed-boundary-holographic-realisation
%  precise: the move from carrier (C3) Z_ch^der(A) to a strict
%  single-colour E_3^top algebra requires a conformal-vector null-
%  homotopy T = [Q, G] in a filtered ambient, with the five obligations
%  below.  Without (T1)-(T5), \SCchtop is the native output and Dunn
%  additivity does NOT apply.
% =====================================================================

\begin{theorem}[Conditional topologisation by stress-tensor null-homotopy;
\ClaimStatusConditional on the listed obligations;
licensing tags $\alpha + \delta + \varepsilon$]%
\label{thm:topologisation-by-stress-tensor-null-homotopy}%
\index{topologisation theorem!conditional on T=[Q,G]|textbf}%
\index{stress tensor null-homotopy!conditional topologisation}%
\index{SCchtop@$\SCchtop$ to $E_{3}^{\mathrm{top}}$ refinement!conditional}%

\noindent\textbf{Status.}\enspace \ClaimStatusConditional on the five
obligations \textup{(T1)}--\textup{(T5)} below.  Without them, the
output is the native $\SCchtop$ bulk-boundary datum of
\cref{thm:chiral-higher-deligne}\,(2) and Dunn additivity does
\emph{not} apply.

\noindent\textbf{Ambient.}\enspace a complete filtered $\Bbbk$-linear
dg category $\mathsf{C}_{\rho}$: for class~$\mathsf{L}$ admitting
\textup{(T3a)}, the original complex $\mathrm{Ch}(\mathrm{Vect})$;
for class~$\mathsf{M}$ admitting only \textup{(T3b)}, the
weight-completed / pro / Banach ambient (hypothesis
$\hypAmbientWtCpl$, $\hypAmbientPro$).

\noindent\textbf{Source.}\enspace a chirally Koszul $E_{1}$-chiral
algebra $\cA$ on a smooth curve $X$ with conformal vector $T(z) \in
\cA(\!(z)\!)$ at non-critical level.

\noindent\textbf{Target.}\enspace a strict $E_{3}^{\mathrm{top}}$
algebra structure on $\Zderch{\cA}$ in $\mathsf{C}_{\rho}$.

\noindent\textbf{Map.}\enspace $\Zderch{\cA} \rightsquigarrow
E_{3}^{\mathrm{top}}\text{-}\mathsf{Alg}\bigl(\mathsf{C}_{\rho}\bigr)$.

\noindent\textbf{Obligations.}\enspace
\begin{description}
\item[\textup{(T1)}] there exist a BRST differential $Q$ on
$\Zderch{\cA}$ and a degree-$(-1)$ operator $G(z) \in \cA(\!(z)\!)$
with
\begin{equation}\label{eq:T-equals-Q-G-bracket}
 T(z) \;=\; \bigl[Q,\; G(z)\bigr]
\end{equation}
at chain level in $\mathsf{C}_{\rho}$, expressing the stress tensor
as a $Q$-coboundary (equivalently, $L_{-1} = \mathrm{Res}_{z}\, T(z)
= [Q,\, \mathrm{Res}_{z}\, G(z)]$, the chiral translation generator
is $Q$-exact);
\item[\textup{(T2)}] $(T, Q, G)$ are compatible with the filtration of
$\mathsf{C}_{\rho}$: each $F_{p}\mathsf{C}_{\rho}$ is preserved by
$Q$ and $G$, $\mathrm{gr}^{F}T \in \mathrm{Im}\,\mathrm{gr}^{F}[Q,-]$,
and the spectral sequence of $F$ converges in $\mathsf{C}_{\rho}$;
\item[\textup{(T3)}] one of
\textup{(a)} \emph{finite propagation in
$\mathrm{Ch}(\mathrm{Vect})$}: $\cA$ has finite-dimensional weight
spaces and the OPE truncates after finitely many singular terms
\textup{(}class~$\mathsf{L}$: affine $\widehat{\fg}_{k}$ at
non-critical $k \ne -h^{\vee}$\textup{)};
\textup{(b)} \emph{weight-completed ambient}: $\mathsf{C}_{\rho}$ is
the pro-object / $J$-adic / Banach completion of
$\mathrm{Ch}(\mathrm{Vect})$ along the conformal-weight filtration,
in which the convolution complex of $\cA$ converges
\textup{(}class~$\mathsf{M}$: $\mathrm{Vir}_{c}$, principal $\cW_{N}$,
hook-type $\cW_{k}(\fg, f)$\textup);
\item[\textup{(T4)}] the brace operations of
\cref{thm:chain-level-bulk}\,(i), transferred along the strong
deformation retract $(i,\,p,\,h)$ associated with
\eqref{eq:T-equals-Q-G-bracket} by the homological perturbation
lemma, are independent of the choice of SDR up to chain homotopy
(hypothesis $\hypKZSDR$ in the analytic case);
\item[\textup{(T5)}] the anomaly classes in
$H^{1}\bigl(\mathrm{C}^{\bullet}_{\mathrm{conv}}\bigr)$ and
$H^{2}\bigl(\mathrm{C}^{\bullet}_{\mathrm{conv}}\bigr)$ of the
convolution complex of the transferred structure vanish (hypothesis
$\hypStokes$ + $\hypReflWts$ in the analytic case;
quasi-classical Maurer--Cartan compatibility in the algebraic case).
\end{description}

\noindent\textbf{Conclusion.}\enspace under \textup{(T1)}--\textup{(T5)}
in $\mathsf{C}_{\rho}$, the mixed $\SCchtop$ datum on
$(\Zderch{\cA}, \cA)$ refines to a strict $E_{3}^{\mathrm{top}}$
algebra structure on $\Zderch{\cA}$,
\begin{equation}\label{eq:topologisation-Dunn-additivity}
 \Zderch{\cA} \;\in\;
 E_{3}^{\mathrm{top}}\text{-}\mathsf{Alg}\bigl(\mathsf{C}_{\rho}\bigr)
 \;=\; E_{2}^{\mathrm{top}}\text{-}\mathsf{Alg}
 \bigl(E_{1}^{\mathrm{top}}\text{-}\mathsf{Alg}(\mathsf{C}_{\rho})\bigr),
\end{equation}
the right-hand identification by Dunn additivity, which applies only
after \textup{(T1)} has rendered the chiral direction $Q$-exact and
\textup{(T3)} has secured the chain-level limit.
\end{theorem}

\begin{proof}
Under \eqref{eq:T-equals-Q-G-bracket} of \textup{(T1)}, the chiral
translation generator $L_{-1} = \mathrm{Res}_{z}\, T(z)$ is the
$Q$-coboundary of $\mathrm{Res}_{z}\, G(z)$; passage to $Q$-cohomology
renders translations along the curve direction trivial.  The
chiral colour of $\SCchtop$, equipped with the chain-level Virasoro
action on $\Zderch{\cA}$ via the brace operations of
\cref{thm:chain-level-bulk}\,(i), descends through $Q$-cohomology
to an $E_{1}^{\mathrm{top}}$ algebra structure on
$H^{\bullet}_{Q}\bigl(\Zderch{\cA}\bigr)$ once \textup{(T2)} ensures
filtration-compatibility of the cohomology passage.

Under \textup{(T3)}, the chain-level limit of the topologised
structure exists in $\mathsf{C}_{\rho}$: in case~\textup{(a)}, by
finite-dimensional weight-space truncation
(\cref{thm:intro-e3-km}); in case~\textup{(b)}, by convergence of
the weight-completion / pro / Banach ambient
(\cref{cor:universal-holography-class-M}).

Under \textup{(T4)}, the homological perturbation lemma associated
with the SDR $(i, p, h)$ of \eqref{eq:T-equals-Q-G-bracket} transfers
the chiral brace operations of \cref{thm:chain-level-bulk}\,(i)
along the projection $p$ and section $i$ to chain-homotopic
operations on the $Q$-cohomology, with the transferred $A_{\infty}$
structure independent of $(i, p, h)$ up to homotopy
(Crainic; Lambe--Stasheff).  Together with the now-topologised chiral
direction this furnishes an $E_{2}^{\mathrm{top}}$-algebra on
$H^{\bullet}_{Q}\bigl(\Zderch{\cA}\bigr)$.

Under \textup{(T5)}, the convolution-complex anomaly classes
vanish in $H^{1}$ and $H^{2}$, so the transferred operations satisfy
the $E_{2}^{\mathrm{top}}$ associativity / brace coherence equations
at chain level without obstruction.

Combining the now-topologised chiral $E_{1}^{\mathrm{top}}$ direction
with the transferred $E_{2}^{\mathrm{top}}$ brace action, Dunn
additivity in $\mathsf{C}_{\rho}$
(Lurie, \emph{Higher Algebra}, Theorem~5.1.2.2; Dunn 1988) yields
\eqref{eq:topologisation-Dunn-additivity}: a strict
$E_{3}^{\mathrm{top}}$-algebra on $\Zderch{\cA}$.

Dunn additivity is invoked \emph{only after} \textup{(T1)} has
topologised the chiral colour and \textup{(T3)} has secured the
chain-level limit; without \textup{(T1)}, the chiral direction is
holomorphic, not topological, and the directional restriction
$\SCchtop(\ldots, \mathsf{top}, \ldots; \mathsf{cl}) = \varnothing$
prevents the Dunn identification (cf.\
Convention~\ref{conv:chd-e3-chiral-meaning}).
\end{proof}

\begin{remark}[Topologisation as residual~(R3) of the typed
boundary-holographic functor]%
\label{rem:topologisation-as-R3-of-Phihol}%
\index{topologisation theorem!as residual (R3) of $\Phihol$}%
The theorem makes residual~\textup{(R3)} of the typed boundary-
holographic functor
(\cref{thm:typed-boundary-holographic-realisation}) precise: the passage from the (C3)
carrier $\Zderch{\cA}$ to the strict $E_{3}^{\mathrm{top}}$ algebra
\eqref{eq:topologisation-Dunn-additivity} is licensed by the data
$(Q, G)$ with $T = [Q, G]$ in the filtered ambient
$\mathsf{C}_{\rho}$, subject to the five obligations.  Without those
data, $\Zderch{\cA}$ carries only the $\SCchtop$-action of
\cref{thm:chiral-higher-deligne}\,(2) and the DS--Hochschild mixed
lift of \cref{thm:chd-ds-hochschild}; the $E_{3}^{\mathrm{top}}$
refinement is strictly further structure, not the native output.
\end{remark}

\begin{remark}[Class~$\mathsf{L}$ versus class~$\mathsf{M}$ ambient
under topologisation]%
\label{rem:class-L-vs-class-M-topologisation}%
\index{class~L versus class~M!topologisation ambient}%
For class~$\mathsf{L}$ (affine $\widehat{\fg}_{k}$ at non-critical
level, with $T(z)$ the Sugawara stress tensor), hypothesis
\textup{(T3a)} holds: chain-level finite propagation gives the
topologisation in $\mathrm{Ch}(\mathrm{Vect})$ directly
(\cref{thm:intro-e3-km}; the BRST identity
$T_{\mathrm{DS}} = [Q_{\mathrm{DS}}, G'_{\mathrm{DS}}]$ of
\cref{thm:intro-e3-ds} is the DS analogue after the total DS
differential and the \(\hypDSBRST\) primitive have been installed). For
class~$\mathsf{M}$ (Virasoro $\mathrm{Vir}_{c}$, principal $\cW_{N}$,
hook-type $\cW_{k}(\fg, f)$), only \textup{(T3b)} holds: the
topologisation is in the weight-completed / pro / Banach ambient,
and the chain-level statement in $\mathrm{Ch}(\mathrm{Vect})$ is
genuinely false in general --- the $\mathrm{Vir}_{c}$ shadow
coefficient $S_{4}(\mathrm{Vir}_{c}) = 10/[c(5c+22)] \ne 0$
populates bar-weight~$4$ and obstructs the
$\mathrm{Ch}(\mathrm{Vect})$ chain-level identification
(\cref{cor:universal-holography-class-M}).  Free PVAs reach
\textup{(T3b)} through the KZ analytic SDR + image-choice Stokes
package $\hypKZSDR + \hypStokes + \hypReflWts + \hypTLift$ of
\cref{thm:intro-e3-free-PVA}.
\end{remark}

\begin{remark}[The false Dunn collapse is the negation of the obligation list]%
\label{rem:forbidden-slogan-3-as-T-violation}%
\index{voice table!slogan~\#3 negation}%
Writing ``$E_{1}$-bar explains $2d \to 3d$'' or ``$\SCchtop = E_{3}$
by Dunn additivity'' collapses the holomorphic colour before
topologisation.  The correct form is the conditional
\cref{thm:topologisation-by-stress-tensor-null-homotopy}: Dunn
additivity is invoked \emph{only after} the topologisation data
$(T, Q, G)$ in the stated filtered ambient have been supplied
and the obligations \textup{(T1)}--\textup{(T5)} verified.  The
$E_{3}^{\mathrm{top}}$ lift is a further structure on the closed
colour, not an operadic identity of $\SCchtop$.
\end{remark}

%%%%%%%%%%%%%%%%%%%%%%%%%%%%%%%%%%%%%%%%%%%%%%%%%%%%%%%%%%%%%%%%%%%%%%%%%%%%%
\section{Filtered $E_3$-PBW transport of first-page support}
\label{sec:chd-rigidity}
%%%%%%%%%%%%%%%%%%%%%%%%%%%%%%%%%%%%%%%%%%%%%%%%%%%%%%%%%%%%%%%%%%%%%%%%%%%%%

Volume~I Theorem~H starts from a family datum $H_H(\cA;S)$ and
identifies $\ChirHoch^\bullet(\cA,\cA)$ with a complex supported in the
finite set~$S$. The route studied here transports an independently
computed first-page support bound through a chiral-$E_3$-PBW spectral
sequence. The PBW presentation supplies the filtered derived centre;
the first-page calculation supplies~$S$.

\begin{conjecture}[Chiral-$E_3$-PBW and first-page support package for
the derived centre;
\ClaimStatusConjectured; licensing tags $\gamma+\delta$]%
\label{conj:chd-e3-pbw-package}%
\index{E_3!PBW package for the derived centre}%
\index{derived chiral centre!PBW package}%
Fix a finite set $S\subset\mathbb Z$.  Let $\cA$ be an $E_1$-chiral
algebra on $X$ in the Vol.~I Koszul locus and let $D_x$ be a formal
disc. The spectral-sequence route requires the following filtered
package.
There must exist:
\begin{enumerate}[label=\textup{(\roman*)}]
\item a complete, separated, exhaustive increasing filtration
$F_{-1}=0\subset F_0\subset F_1\subset\cdots$ on
\[
 R_x := Z^{\mathrm{der}}_{\mathrm{ch}}(\cA|_{D_x})
      = \ChirHoch^\bullet(\cA|_{D_x},\cA|_{D_x})
\]
by arity / collision weight, compatible with the topologised
$E_3$-operations in the sense that every $r$-ary operation sends
$F_{p_1}\otimes\cdots\otimes F_{p_r}$ into
$F_{p_1+\cdots+p_r}$, with finite-dimensional graded pieces in each
fixed conformal weight;
\item a local graded dg Lie object \(W_x\) with a specified tangent
identification
\[
 W_x \simeq
 \mathbb T^{E_3}_{R_x,x}[-2]
\]
in the completed filtered ambient, and a filtered $E_3$-algebra
quasi-isomorphism
\begin{equation}\label{eq:chd-e3-pbw-filtered}
 R_x \;\simeq\; U_{E_3}(W_x);
\end{equation}
\item an associated-graded $E_3$-algebra isomorphism
\begin{equation}\label{eq:chd-e3-pbw-gr}
 \theta_x\colon
 \Free_{E_3}\!\bigl(H^\bullet(W_x[-2])\bigr)
 \xrightarrow{\;\sim\;}
 \operatorname{gr}_{F} R_x,
\end{equation}
equivalently \(\ker\theta_x=\operatorname{coker}\theta_x=0\). Thus
the operations encoded by~\(W_x\) generate the associated graded with
precisely their divided-power, Orlik--Solomon, and collision-residue
relations.
Moreover \(W_x\) must be the primitive source of this associated
graded: the induced map to \(E_3\)-indecomposables is an isomorphism
\begin{equation}\label{eq:chd-e3-pbw-primitives}
 H^\bullet(W_x[-2])
 \xrightarrow{\;\sim\;}
 Q_{E_3}\!\left(\operatorname{gr}_{F}R_x\right)
 :=
 \operatorname{gr}_{F}R_x\big/
 \bigl(E_3\text{-decomposable operations}\bigr).
\end{equation}
Thus \(W_x\) is the primitive tangent datum whose free
\(E_3\)-envelope gives the associated graded;
\item a Rees-flat lift. With
\[
 \mathcal R_F R_x:=\bigoplus_{p\ge0}F_pR_x\,u^p
\]
one has
\begin{equation}\label{eq:chd-e3-pbw-rees}
 \mathcal R_F R_x/(u)\cong \operatorname{gr}_F R_x,\qquad
 \mathcal R_F R_x[u^{-1}]_{u=1}\simeq R_x,
 \qquad H^\bullet(\mathcal R_F R_x)\text{ is }u\text{-torsion-free};
\end{equation}
\item the strongly convergent PBW spectral sequence
\begin{equation}\label{eq:chd-e3-pbw-ss}
 E_1^{p,q}
 =
 H^{p+q}\!\left(\operatorname{gr}_F^pR_x\right)
 \cong
 H^{p+q}\!\left(\Free^p_{E_3}(W_x[-2])\right)
 \;\Longrightarrow\;
 H^{p+q}(R_x);
\end{equation}
\item polynomial growth and amplitude bounds: for the conformal-weight
grading $n$ on $H^q(W_x)$ there are constants $C,d,N_0$ such that
\begin{equation}\label{eq:chd-e3-pbw-growth}
 \dim H^q(W_x)_n \le C(1+n)^d,\qquad
 H^q(W_x)=0\quad(q>N_0),
\end{equation}
and the shifted $E_1$ page satisfies the support bound
\begin{equation}\label{eq:chd-e3-pbw-page-bound}
 E_1^{p,q}=0 \qquad (p+q\notin S).
\end{equation}
\end{enumerate}
The established Theorem~H proof uses the deformation retract contained
in $H_H(\cA;S)$. The package above transports an independent
first-page support calculation through the $E_3$-PBW filtration.
\end{conjecture}

\begin{lemma}[Spectral-sequence transport of first-page support]%
\label{lem:chd-e3-rigidity-point}%
\index{E_3!PBW support transport}%
If Conjecture~\ref{conj:chd-e3-pbw-package} holds at $x$, then
\[
 \operatorname{Supp}H^\bullet
 Z^{\mathrm{der}}_{\mathrm{ch}}(\cA|_{D_x})\subseteq S.
\]
\end{lemma}

\begin{proof}
Apply the spectral sequence~\eqref{eq:chd-e3-pbw-ss}. The filtration is
complete, separated, exhaustive, and bounded below in arity, while
\eqref{eq:chd-e3-pbw-rees} supplies the $u$-torsion-free Rees lift;
strong convergence is part of the PBW package. The support bound
\eqref{eq:chd-e3-pbw-page-bound} gives \(E_1^{p,q}=0\) whenever
\(p+q\notin S\).  Each subsequent term \(E_{r+1}^{p,q}\) is a
subquotient of \(E_r^{p,q}\), so the same support condition holds on
every page and on the abutment. The polynomial-growth estimate
\eqref{eq:chd-e3-pbw-growth} supplies the finiteness needed for the
filtered $E_3$ envelope and the local-to-global passage; the
first-page calculation supplies the support bound.
\end{proof}

\begin{remark}[The page-support clause is independent]
\label{rem:chd-e3-pbw-page-support-independent}
The support bound \eqref{eq:chd-e3-pbw-page-bound} is a separate
vanishing assertion.  Polynomial growth controls dimensions and leaves
the support set undetermined.  For example, take
$S_0=\{0,1\}$ and consider a first-quadrant page with
\[
 E_1^{0,0}=\Bbbk,\qquad
 E_1^{1,0}=\Bbbk,\qquad
 E_1^{3,0}=\Bbbk,
 \qquad
 E_1^{p,q}=0\ \text{otherwise}.
\]
All dimensions are polynomially bounded and the amplitude is finite.
The cohomological differential on the PBW page has bidegree
\((r,1-r)\), so it raises total degree by \(1\). The neighboring groups
in total degrees~\(2\) and~\(4\) vanish, hence the class in
\(E_1^{3,0}\) survives to the abutment of every strongly convergent
spectral sequence with this page. Thus a
polynomial-growth \(E_3\)-envelope can still carry degree-\(3\)
cohomology. The independent input in
Conjecture~\ref{conj:chd-e3-pbw-package} is the support statement
\(E_1^{p,q}=0\) for \(p+q\notin S\). Finite generation, polynomial
growth, and bounded conformal weight are separate conditions.
\end{remark}

\begin{conjecture}[Family support transported through filtered
chiral-$E_3$-PBW]%
\label{conj:H-concentration-via-E3-rigidity}%
\index{Theorem H!consequence of E_3 rigidity}%
Fix a finite set $S\subset\mathbb Z$.  Let $\cA$ be an $E_1$-chiral
algebra on $X$ in the Koszul locus with Vol.~I PBW collapse. Suppose
the derived centres on formal discs admit
the chiral-$E_3$-PBW package of
Conjecture~\ref{conj:chd-e3-pbw-package}, compatibly with
Gelfand--Fuchs--Beilinson--Drinfeld local-to-global. Then the
first-page support clause, strong convergence, and support-preserving
descent (Lemma~\ref{lem:chd-e3-rigidity-point}, applied pointwise to
$Z^{\mathrm{der}}_{\mathrm{ch}}(\cA|_{D_x})$) give
\[
 \operatorname{Supp}\ChirHoch^\bullet(\cA,\cA)\subseteq S.
\]
The required input is the Rees-flat associated-graded and growth package
\eqref{eq:chd-e3-pbw-filtered}--\eqref{eq:chd-e3-pbw-page-bound} (see
Remark~\ref{rem:H-concentration-via-E3-rigidity-type-error-retraction}).
Volume~I Theorem~H proves this conclusion from the independent family
datum $H_H(\cA;S)$. The filtered route transports the same support
from the independently computed PBW page.
\end{conjecture}

\begin{remark}[Typed comparison required by the filtered $E_3$ route]%
\label{rem:H-concentration-via-E3-rigidity-type-error-retraction}%
\index{Theorem H!E_3-rigidity route obstruction}%
\index{native vs derived E_3!type error}%
Theorem~\ref{thm:chiral-higher-deligne} places the topologised
$E_3$-structure on
$Z^{\mathrm{der}}_{\mathrm{ch}}(\cA|_{D_x})$.  Volume~I PBW collapse
places an $E_1$-chiral envelope structure on $\cA|_{D_x}$. The
filtered argument therefore requires a third comparison: the
$E_3$-presentation
\[
 Z^{\mathrm{der}}_{\mathrm{ch}}(\cA|_{D_x})
 \simeq U_{E_3}(W_x)
\]
with associated graded, Rees-flat lift, strongly convergent spectral
sequence, and first-page support in~$S$.  This is precisely
Conjecture~\ref{conj:chd-e3-pbw-package}.  Volume~I Theorem~H follows
by the independent deformation retract in $H_H(\cA;S)$.
\end{remark}

\begin{remark}[Proof route under the filtered transport
Conjecture~\ref{conj:H-concentration-via-E3-rigidity}]%
\label{rem:H-concentration-via-E3-rigidity-proof-route}%
\index{Theorem H!E_3-rigidity proof route}%
Fix the finite set~$S$ in
Conjecture~\ref{conj:chd-e3-pbw-package}.  On a formal disc~$D_x$,
Theorem~\ref{thm:chiral-higher-deligne} gives the topologised
$E_3$-structure on
$R_x=Z^{\mathrm{der}}_{\mathrm{ch}}(\cA|_{D_x})$.  The conjectural
filtered presentation gives the spectral sequence
\eqref{eq:chd-e3-pbw-ss}, whose first page is supported in~$S$.
Lemma~\ref{lem:chd-e3-rigidity-point} carries this support to
$H^\bullet(R_x)$.  A support-preserving
Gelfand--Fuchs--Beilinson--Drinfeld descent then gives
\[
 \operatorname{Supp}\ChirHoch^\bullet(\cA,\cA)\subseteq S.
\]
Thus the route consists of the topologised $E_3$ structure, the
filtered PBW presentation, and support-preserving descent.  Volume~I
Theorem~H reaches the same conclusion through the deformation retract
in $H_H(\cA;S)$.
\end{remark}

\begin{remark}[Family support as a consequence of chiral-$E_3$-PBW]%
\label{rem:chd-consequence-not-hypothesis}%
\index{Theorem H!consequence vs hypothesis}%
The derived-centre route depends on the chiral-$E_3$-PBW theorem
of Step~2: the filtered associated-graded identification, Rees-flat
lift, convergent spectral sequence, polynomial-growth/amplitude bounds,
and support bound of Conjecture~\ref{conj:chd-e3-pbw-package}.  Such a
presentation computes the family support~$S$ from the first page.
The canonical Volume~I route computes the same support from the
explicit maps and homotopy in $H_H(\cA;S)$.
\end{remark}

%%%%%%%%%%%%%%%%%%%%%%%%%%%%%%%%%%%%%%%%%%%%%%%%%%%%%%%%%%%%%%%%%%%%%%%%%%%%%
\section{DS-Hochschild compatibility at chain level}
\label{sec:chd-ds}
%%%%%%%%%%%%%%%%%%%%%%%%%%%%%%%%%%%%%%%%%%%%%%%%%%%%%%%%%%%%%%%%%%%%%%%%%%%%%

The Drinfeld-Sokolov reduction $V_k(\fg) \rightsquigarrow \cW_k(\fg, f)$
(principal or hook-type) is the passage from affine Kac-Moody to
$\cW$-algebras through BRST reduction at a nilpotent orbit. The
principal/hook-type DS deformation retract gives a chain-level
comparison between chiral Hochschild complexes.

Fix a good grading
\[
\fg=\bigoplus_{j\in \frac1q\mathbb Z}\fg_j
\]
for \(f\), and choose a homogeneous basis
\(\{u_\alpha\}\) of \(\fg_{>0}\), with ghosts
\(\psi_\alpha,\psi_\alpha^*\).  The quantum DS algebra is the degree
zero cohomology
\begin{equation}
\label{eq:ds-brst-w-algebra}
\cW_k(\fg,f)
=
H^0\!\left(V_k(\fg)\widehat\otimes\cF_{\mathrm{gh}},Q_{\mathrm{DS}}\right),
\end{equation}
where the BRST charge has KRW form
\begin{equation}
\label{eq:ds-brst-charge}
Q_{\mathrm{DS}}
=
\operatorname*{Res}_{z=0}
\left[
\sum_{\alpha}
\bigl(J^{u_\alpha}(z)-\chi(u_\alpha)\bigr)\psi_\alpha^*(z)
-\frac12
\sum_{\alpha,\beta,\gamma}
f_{\alpha\beta}^{\gamma}
:\psi_\alpha^*(z)\psi_\beta^*(z)\psi_\gamma(z):
+Q_{\mathrm{neut}}(z)
\right]\,dz .
\end{equation}
Here \(Q_{\mathrm{neut}}\) is the standard neutral-fermion term attached
to the half-integer part of the good grading; it vanishes in the
integral principal case.  The Hochschild comparison is required at the
chain level to be a special deformation retract
\begin{equation}
\label{eq:ds-hochschild-sdr}
\left(
\ChirHoch^\bullet(V_k(\fg),V_k(\fg))\widehat\otimes
\cF_{\mathrm{gh}},Q_{\mathrm{DS}}
\right)
\;\substack{\xrightarrow{\;p\;}\\[-2pt]\xleftarrow[\;i\;]{} }\;
\ChirHoch^\bullet(\cW_k(\fg,f),\cW_k(\fg,f)),
\qquad
Q_{\mathrm{DS}}h+hQ_{\mathrm{DS}}=\mathrm{id}-i\circ p .
\end{equation}

\begin{theorem}[DS-Hochschild chain-level compatibility with mixed lift;
\ClaimStatusConditional; licensing tags $\alpha+\beta+\gamma+\delta$ via
good grading, DS SDR, $\hypAmbientWtCpl$, and mixed lift hypotheses]%
\label{thm:chd-ds-hochschild}%
\index{Drinfeld-Sokolov!Hochschild compatibility}%
Let $\fg$ be a simple Lie algebra, $f \in \fg$ a nilpotent of principal
or hook type, and $k \ne -h^{\vee}$. Write $\cW_k(\fg, f)$ for the
corresponding $\cW$-algebra at level $k$. Assume:
\begin{enumerate}[label=\textup{(\alph*)}]
\item a good grading has been fixed; in fractional degree the cover
 \(z=s^q\) and the \(\mu_q\)-descent datum of
 Remark~\ref{rem:ds-fractional-good-grading-cover} are supplied;
\item the DS-Hochschild complex is taken in the completed ambient
 \(\hypAmbientWtCpl\) and admits the special deformation retract
 \eqref{eq:ds-hochschild-sdr}, induced from a completed
 DS bar-coalgebra SDR on the reduced chiral bar coalgebras, with
 \(p i=\mathrm{id}\), \(ph=hi=h^2=0\), and
 \(Q_{\mathrm{DS}}h+hQ_{\mathrm{DS}}=\mathrm{id}-i\circ p\);
\item the chiral Hochschild complex is represented by the completed
 coderivation model
 \(\Coder_0(B^{\mathrm{ch}}(-))\), the bar-coalgebra SDR in
 \textup{(b)} induces the displayed coderivation SDR by transfer, and
 the HPL tree sums satisfy the bounded-shift/weightwise-convergence
 estimates in \(\hypAmbientWtCpl\);
\item for the mixed chiral-topological refinement, the heptagon
 hypotheses and two-coloured cobar homotopy of
 Remark~\ref{rem:chd-universality-hypotheses} are supplied.
\end{enumerate}
Then at chain level there is
an $\infty$-quasi-isomorphism of $E_2$-chiral brace algebras (equivalently,
a quasi-iso on cohomology of the strict Gerstenhaber structures)
\begin{equation}\label{eq:chd-ds-hh}
 \mathrm{DS}^\bullet:
 \ChirHoch^\bullet(\cW_k(\fg, f)) \;\xrightarrow{\;\sim\;}\;
 H^\bullet_{\mathrm{DS}}\bigl(\ChirHoch^\bullet(V_k(\fg))\bigr),
\end{equation}
where $H^\bullet_{\mathrm{DS}}$ denotes the DS cohomology of the
BRST complex with nilpotent reduction along $f$. The quasi-isomorphism
is the HPL transfer along~\eqref{eq:ds-hochschild-sdr}: for every
planar rooted tree \(T\) with \(n\) leaves the \(n\)-ary brace is
\begin{equation}
\label{eq:ds-hochschild-hpl-brace}
\mu_{\cW}^{n}
=
\sum_{T\in\mathrm{PRT}_n}
\pm\;
p\,\mu^V_T(i,\ldots,i;h,\ldots,h).
\end{equation}
Equivalently,
\[
\ChirHoch^\bullet(\cW_k(\fg,f),\cW_k(\fg,f))
\simeq
H^\bullet_{\mathrm{DS}}\!
\left(\ChirHoch^\bullet(V_k(\fg),V_k(\fg))\right)
\]
as \(E_2^{\mathrm{ch}}\)-brace algebras.  The quasi-isomorphism
is compatible with the mixed chiral-topological extensions of
Thm~\ref{thm:chiral-higher-deligne}
(Convention~\ref{conv:chd-e3-chiral-meaning}) after the heptagon
hypotheses, the two-coloured cobar homotopy of
Remark~\ref{rem:chd-universality-hypotheses}, and the completed mixed
analytic estimates.
\end{theorem}

\begin{proof}
Four steps: DS complex, chiral HKR model, HPL, intertwining.

Step 1: DS complex and retract. Kac--Roan--Wakimoto construct the
quantum Hamiltonian-reduction BRST complex
\((V_k(\fg)\widehat\otimes\cF_{\mathrm{gh}},Q_{\mathrm{DS}})\).
Arakawa's representation-theoretic results identify the relevant
DS cohomology on the stated principal/hook-type locus after the
standard finiteness hypotheses are imposed.  The theorem does not use
the false assertion that every non-critical generic level is
\(C_2\)-cofinite; the required input is exactly the special
deformation retract in hypothesis \textup{(b)}.

Step 2: chiral HKR/brace model for each side. The chosen chiral
Hochschild model of hypothesis \textup{(c)}, equivalently the
polyvector presentation on the locus where the chiral HKR comparison is
available, gives
\[
 \ChirHoch^\bullet(\cW_k(\fg, f)) \;\simeq\;
 \Gamma\bigl(\operatorname{Spec}(\cW_k^\flat), \bigwedge^\bullet \cT\bigr)^{\cW}_{\mathrm{chiral}},
\]
where $\cT$ is the tangent complex and the right-hand side uses the
chiral polyvector sheaf. Analogously for $V_k(\fg)$ on the other side
of \eqref{eq:chd-ds-hh}.

Step 3: HPL transfer through the DS coderivation retract. Hochschild
cochains are not functorial in an algebra map, so
\eqref{eq:ds-hochschild-sdr} is not obtained by applying
\(\ChirHoch^\bullet(-,-)\) to the DS algebra retract.  The data in
hypotheses \textup{(b)}--\textup{(c)} are stronger: a completed
DS homotopy equivalence of reduced chiral bar coalgebras
\[
\barBch(\cW_k(\fg,f))
\;\rightleftarrows\;
\barBch(V_k(\fg))\widehat\otimes\cF_{\mathrm{gh}}
\]
with homotopy \(h_B\), and the induced transfer on completed
coderivation complexes
\[
\Coder_0(\barBch(\cW_k(\fg,f)))
\;\rightleftarrows\;
\Coder_0(\barBch(V_k(\fg))\widehat\otimes\cF_{\mathrm{gh}}).
\]
Conjugation by the transferred bar-coalgebra maps and the standard
coderivation HPL give the displayed Hochschild SDR.  The side
conditions \(p\circ i=\mathrm{id}\), \(ph=hi=h^2=0\), and
\(Q_{\mathrm{DS}}h+hQ_{\mathrm{DS}}=\mathrm{id}-i\circ p\) make the
homological-perturbation theorem applicable in the completed chiral
brace category.

The HPL transfer formula is precisely
\eqref{eq:ds-hochschild-hpl-brace}: internal vertices are decorated by
the affine Hochschild braces \(\mu^V\), leaves by \(i\), the root by
\(p\), and internal edges by \(h\).  Finiteness of the HPL sum does
not come from generic \(C_2\)-cofiniteness.  It comes from the Kazhdan
and conformal-weight filtrations on the BRST complex, the finite
dimensionality of the relevant graded pieces, and the bounded-shift
estimate in hypothesis \textup{(c)}.  On the original direct-sum
complex one still needs the corresponding Banach or bounded-shift
estimate; without it the theorem asserts no chain-level direct-sum
comparison.

Step 4: DS-bar intertwining. The intertwining with bar complexes
(Vol.~I, Theorem~\ref*{V1-thm:ds-koszul-intertwine}):
$\Barch(\cW_k(\fg, f)) \simeq H^\bullet_{\mathrm{DS}}(\Barch(V_k(\fg)))$,
is the source of the bar-coalgebra homotopy equivalence used in
Step~3.  The coderivation transfer along this equivalence promotes
the Step-3 $\infty$-quasi-isomorphism of brace algebras to a
quasi-isomorphism of $E_2$-chiral brace algebras (and of strict
Gerstenhaber algebras on cohomology). The mixed chiral-topological lift
to a quasi-isomorphism of $\SCchtop$ bulk-boundary data is obtained via
Part~(2) of Thm~\ref{thm:chiral-higher-deligne}: the heptagon
hypotheses supply the two-coloured target and
Remark~\ref{rem:chd-universality-hypotheses} supplies the exact
cobar-homotopy obligation. On cohomology the associator dependence
collapses by Part~(3), so the Gerstenhaber comparison needs no
two-coloured cobar homotopy. The $E_2$-chiral brace-algebra
quasi-isomorphism is the chain-level content of the theorem.

Conclusion. The chain-level compatibility \eqref{eq:chd-ds-hh} is the
DS-Hochschild bridge used by the class-$M$ global triangle,
holographic reconstruction, symplectic polarization, and universal
holography constructions, each through the same explicit chain map.
\end{proof}

\begin{remark}[Subregular and minimal nilpotents]%
\label{rem:chd-nonprincipal-nilpotents}%
The hypothesis on $f$ restricts Theorem~\ref{thm:chd-ds-hochschild} to
principal and hook-type nilpotents. For subregular and minimal $f$,
the coset presentation of $\cW_k(\fg,f)$ (Kac--Roan--Wakimoto,
\emph{Quantum reduction for affine superalgebras}, Comm.~Math.~Phys.\
2003; Arakawa--Creutzig--Linshaw coset-duality) exhibits
$\cW_k(\fg,f)$ as a commutant of a smaller affine $\cW$ times an
auxiliary free algebra; the DS--Hochschild bridge extends through the
coset commutant, and the details, parallel to the principal case,
appear in the $\cW$-landscape treatment of Vol.~I. Exotic nilpotents
remain open.
\end{remark}

\begin{remark}[Fractional good gradings and cover descent]
\label{rem:ds-fractional-good-grading-cover}
For a non-principal nilpotent with good grading in
\(\frac1q\mathbb Z\), the DS HPL transfer is performed first on the
branched cover \(z=s^q\).  A field of Kazhdan degree \(j/q\) becomes
integral degree \(j\) on the \(s\)-cover.  The antighost primitive
\(G'_f\) and stress tensor satisfy
\[
T_{\mathrm{DS}}=[Q_{\mathrm{DS}},G'_f]
\]
upstairs.  If the transferred braces and \(G'_f\) are \(\mu_q\)-invariant,
then they descend to the original \(z\)-disc.  This is the mechanism for
Bershadsky--Polyakov and hook-type cases; without the cover descent
datum the non-principal chain-level \(\SCchtop\) lift is not asserted.
\end{remark}

%%%%%%%%%%%%%%%%%%%%%%%%%%%%%%%%%%%%%%%%%%%%%%%%%%%%%%%%%%%%%%%%%%%%%%%%%%%%%
\section{Universal holography at chain level for class M}
\label{sec:chd-holography}
%%%%%%%%%%%%%%%%%%%%%%%%%%%%%%%%%%%%%%%%%%%%%%%%%%%%%%%%%%%%%%%%%%%%%%%%%%%%%

\begin{corollary}[Class-M bulk-centre comparison in the
weight-completed ambient; \ClaimStatusConditional; licensing tags
$\beta+\gamma+\delta$]%
\label{cor:universal-holography-class-M}%
\index{universal holography!class M weight-completed ambient}%
Let $\cA$ be a class-M chiral algebra
\textup{(}Virasoro, principal $\cW_N$, or verified hook-type
$\cW_k(\fg,f)$\textup{)} at generic non-critical level.  Assume:
\begin{enumerate}[label=\textup{(\alph*)}]
\item the class-M bar/BRST complexes are taken in the strict
Mittag--Leffler weight-completed, pro-, or $J$-adic ambient
\(\hypAmbientWtCpl\);
\item the DS--Hochschild transport theorem
\textup{(}Theorem~\ref{thm:chd-ds-hochschild}\textup{)} applies to the
chosen principal or hook-type reduction, including the required cover
descent in fractional good gradings;
\item a 3d HT BV/factorization model has been constructed for
\(\cA\), and its local bulk observables are identified with the
derived chiral centre by a bulk--Hochschild comparison map
\[
\chi_{\mathrm{HT}}\colon
 Z^{\mathrm{der}}_{\mathrm{ch}}(\cA)
 \longrightarrow
 \mathcal A_{\mathrm{bulk}}^{\mathrm{3d\,HT}}(\cA);
\]
\item for the chain-level mixed chiral-topological refinement, the
heptagon hypotheses and the two-coloured cobar homotopy of
Remark~\ref{rem:chd-universality-hypotheses} are supplied.
\end{enumerate}
Then \(\chi_{\mathrm{HT}}\) is an equivalence on cohomology as an
\(E_2\)-chiral/Gerstenhaber algebra,
\[
 Z^{\mathrm{der}}_{\mathrm{ch}}(\cA) \;\simeq\;
 \mathcal{A}_{\mathrm{bulk}}^{\mathrm{3d\,HT}}(\cA),
\]
and, under \textup{(d)}, refines to a mixed chiral-topological
chain-level equivalence in the sense of
Convention~\ref{conv:chd-e3-chiral-meaning}. On the bounded
direct-sum ambient $\mathrm{Ch}(\mathrm{Vect})$ the class-M
chain-level identification is false: already
$S_4(\mathrm{Vir}_c)=10/[c(5c+22)]\ne0$ populates bar-weight~$4$.
The G/L/C classes are governed by their own original-complex
identifications and are not consequences of the class-M comparison.
\end{corollary}

\begin{proof}
The direct-sum complex fails in class M, so the comparison is made in
the ambient of hypothesis \textup{(a)}.  Theorem~\ref{thm:chd-ds-hochschild}
and hypothesis \textup{(b)} identify the class-M chiral Hochschild
complex with the DS cohomology of the corresponding affine chiral
Hochschild complex.  Hypothesis \textup{(c)} is exactly the
field-theoretic input identifying the resulting derived chiral centre
with the local bulk observables of the chosen 3d HT model.  Therefore
\(\chi_{\mathrm{HT}}\) is a cohomological equivalence of
\(E_2\)-chiral/Gerstenhaber algebras.

For the chain-level statement one also needs hypothesis
\textup{(d)}.  Clause~(2) of Theorem~\ref{thm:chiral-higher-deligne}
then supplies the two-coloured \(\SCchtop\) lift, and the mixed-lift
clause of Theorem~\ref{thm:chd-ds-hochschild} transports it through
the DS retract.  These operations converge only in the completed
ambient from \textup{(a)}.  The displayed nonzero Virasoro
\(S_4\)-coefficient gives the obstruction on the bounded direct-sum
ambient.
\end{proof}

\begin{remark}[Consequences for the universal holography part]%
\label{rem:chd-vol2-part-vi}%
When the hypotheses of Corollary~\ref{cor:universal-holography-class-M}
are supplied, the
Universal Holography functor
(Chapter~\ref{ch:thqg-holographic-reconstruction}) has a class-M
weight-completed branch compatible with the G/L/C original-complex
branches.  The
direct-sum class-M complex carries the nonzero class
$S_4(\mathrm{Vir}_c)=10/[c(5c+22)]$ in bar-weight~$4$; this is the
mathematical obstruction forcing the class-M ambient.
\end{remark}

%%%%%%%%%%%%%%%%%%%%%%%%%%%%%%%%%%%%%%%%%%%%%%%%%%%%%%%%%%%%%%%%%%%%%%%%%%%%%
\section{Proof inputs and obstruction classes}
\label{sec:chd-proof-inputs}
%%%%%%%%%%%%%%%%%%%%%%%%%%%%%%%%%%%%%%%%%%%%%%%%%%%%%%%%%%%%%%%%%%%%%%%%%%%%%

\begin{enumerate}
\item \textbf{thm:chiral-higher-deligne} (\S\ref{sec:chd-main}).
The $E_2$-brace action is the chiral Deligne--Tamarkin action of
Thm~\ref{thm:chd-deligne-tamarkin}. The two-coloured lift uses:
(a) chiral-$\SCchtop$ pentagon edges
$(3)\leftrightarrow(4)\leftrightarrow(5)$ of the heptagon
(this volume, Ch.~\ref{ch:sc-chtop-heptagon}); (b) the ordered
bar-cobar action of the two-coloured $\SCchtop$ cooperad on the pair
$(Z^{\mathrm{der}}_{\mathrm{ch}}(\cA),\cA)$. A Dunn theorem is not
used to identify $\SCchtop$ with a single-colour $E_3$ operad.
Tamarkin 2003 (\emph{Another proof of Deligne's conjecture},
Lett.~Math.~Phys.\ 66) supplies the $E_2$ brace theorem at the formal
disc; Kontsevich 2006 (arXiv:math/0608180) supplies the
Stokes/logarithmic-form brace model; Francis 2012
(\emph{The tangent complex and Hochschild cohomology of $E_n$-rings})
supplies the chiral Deligne identification at $E_2$ on a curve,
independent of $\SCchtop$. The missing chain-level datum is
the signed tree homotopy \(h_{\mathrm{oc}}^{\mathrm{tree}}\) and the
vanishing of its obstruction class
\eqref{eq:chd-oc-obstruction-class}.

\item \textbf{conj:H-concentration-via-E3-rigidity}
(\S\ref{sec:chd-rigidity}; route through the derived-centre
$E_3$-PBW obstruction in
Remark~\ref{rem:H-concentration-via-E3-rigidity-type-error-retraction}).
It uses: (a) Lemma~\ref{lem:chd-e3-rigidity-point} via
$E_3$-Koszul duality (Fresse, vol.~II, \S13); (b) chiral PBW
collapse in the Koszul locus
(Vol.~I~\ref*{V1-thm:pbw-koszulness-criterion}); (c) Gel'fand--Fuchs
local-to-global for constructible sheaves of bounded amplitude.
The proof obligation is the chiral-$E_3$-PBW package
of Conjecture~\ref{conj:chd-e3-pbw-package}: filtered presentation,
associated graded, Rees-flat lift, convergent PBW spectral sequence, polynomial
growth, amplitude bound, and \(E_1\)-page support in total degrees
belonging to the chosen set~$S$.  Volume~I reaches the support
conclusion through the explicit retract in $H_H(\cA;S)$.

\item \textbf{thm:chd-ds-hochschild} (\S\ref{sec:chd-ds}).
The $E_2$-chiral brace quasi-isomorphism uses:
(a) the KRW/Arakawa DS BRST complex and its degree-zero cohomology
statement at non-critical level; (b) the completed bar-coalgebra
DS SDR and the induced coderivation transfer
\(\Coder_0(B^{\mathrm{ch}}(-))\); (c) completed HPL with the
bounded-shift/weightwise-convergence estimate. Chiral HKR or
polyvector coordinates may identify the resulting coderivation complex
on smooth loci, but they are not the source of functoriality.
Kac--Wakimoto 1988 (\emph{Modular invariance representations of
affine algebras}) supplies an independent DS cohomology
characterization of $\cW$-characters; Feigin--Frenkel coset duality
for $\cW_k(\fg) \cong \mathrm{coset}$ supplies a coset-commutant
check of the same low-degree character shadow, not a replacement for
the chain-level coderivation SDR. The mixed lift uses the same
two-coloured cobar homotopy
\(h_{\mathrm{oc}}^{\mathrm{tree}}\) of
Definition~\ref{def:chd-two-coloured-cobar-tree-homotopy}.
\end{enumerate}

%%%%%%%%%%%%%%%%%%%%%%%%%%%%%%%%%%%%%%%%%%%%%%%%%%%%%%%%%%%%%%%%%%%%%%%%%%%%%
\section{Proof obligations not discharged here}
\label{sec:chd-not-proved}
%%%%%%%%%%%%%%%%%%%%%%%%%%%%%%%%%%%%%%%%%%%%%%%%%%%%%%%%%%%%%%%%%%%%%%%%%%%%%

\begin{enumerate}
\item Off-locus chain-level extension. Theorem~\ref{thm:chiral-higher-deligne}
is on the Koszul locus. Off-locus the mixed chiral-topological action
extends to the
weight-completed category (Vol.~I~\texttt{thm:mc5-class-m-chain-level-pro-ambient}),
not to the direct-sum chain-level category. The proof at chain level
off-locus is open.

\item Exotic nilpotents in Theorem~\ref{thm:chd-ds-hochschild}. The
principal and hook-type cases are covered; subregular and minimal
reduce via cosets (Remark~\ref{rem:chd-nonprincipal-nilpotents}); exotic
remains open.

\item Chiral-$E_3$-PBW package in
Conjecture~\ref{conj:chd-e3-pbw-package}. Bakalov--De Sole--Kac
\cite{BDSK21} \textup{(}Theorems~7.4 and~7.2\textup{)} give the bounded
benchmarks $(2,1,0,\ldots)$ for the rank-one even superboson and
support $\{0,2,3\}$ for Virasoro; their affine statement retains the
status of Conjecture~7.5. Each chart calculation uses the
corresponding bounded-to-chart quasi-isomorphism.  The $E_3$-PBW route
additionally asks for the filtered derived-centre presentation,
associated graded~\eqref{eq:chd-e3-pbw-gr}, Rees-flat
lift~\eqref{eq:chd-e3-pbw-rees}, convergent spectral
sequence~\eqref{eq:chd-e3-pbw-ss}, and page support
\eqref{eq:chd-e3-pbw-page-bound}.

\item Associator-free chain-level structure. Chain-level associator
dependence is not eliminable; the chain-level mixed
chiral-topological structure, and any strict $E_3^{\mathrm{top}}$
refinement under topologisation, exists only after $\Phi$ is fixed
(Theorem~\ref{thm:chiral-higher-deligne}(3)). Cohomologically the
binary brace structure is associator-free.

\item The MC4 resonance locus. On the resonance locus the DS retract
of Step 3 of Theorem~\ref{thm:chd-ds-hochschild} degenerates
(Feigin-Frenkel kernel collapses at $k = -h^\vee$); the theorem
excludes the critical level by hypothesis.
\end{enumerate}

These are open mathematical questions.

%%%%%%%%%%%%%%%%%%%%%%%%%%%%%%%%%%%%%%%%%%%%%%%%%%%%%%%%%%%%%%%%%%%%%%%%%%%%%
%% End of chapter.
%%%%%%%%%%%%%%%%%%%%%%%%%%%%%%%%%%%%%%%%%%%%%%%%%%%%%%%%%%%%%%%%%%%%%%%%%%%%%
